\documentclass{article}

\PassOptionsToPackage{numbers}{natbib}
\usepackage[final]{neurips_2019}

\usepackage[utf8]{inputenc}
\usepackage[T1]{fontenc}
\usepackage{hyperref}
\usepackage{url}
\usepackage{booktabs}
\usepackage{amsfonts}
\usepackage{amsmath}
\usepackage{microtype}
\usepackage{graphicx}
\usepackage{xcolor}
\usepackage{xparse}

\newif\ifshowinstructions
\showinstructionsfalse

\newcommand{\note}[1]{\ifshowinstructions#1\fi}
\NewDocumentEnvironment{instructions}{+b}
  {\ifshowinstructions #1\par\fi}
  {}

\NewDocumentCommand{\projecttitle}{m m m O{} O{}}{%
  \title{%
    #1\\[0.7em]
    {\normalsize\normalfont
      Columbia COMS 4705 Project Milestone
    }\\[0.8em]
    \parbox{0.95\textwidth}{%
      \normalfont\small\centering

      \textbf{Keywords:} \textit{#2}
      \par\vspace{0.4em}

      \textbf{TA Mentor:} #3%

      \if\relax\detokenize{#4}\relax\else
        \par\vspace{0.4em}
        \textbf{External Collaborators:} #4%
      \fi

      \if\relax\detokenize{#5}\relax\else
        \par\vspace{0.4em}
        \textbf{Sharing Project:} #5%
      \fi
      \par
    }%
  }%
}

\projecttitle
  {What Does a Wish Leave Out? Requirement Recovery in Wish-Driven Code Agents}
  {wish-driven programming, code agents, requirement elicitation, underspecification}
  {[TA name]} % TODO: fill in the mentor who graded the proposal.
  []
  []

\author{
  Yanhao Bai \quad Kaitlyn Bassford \quad Wen Tao Koh \\
  Department of Computer Science, Columbia University \\
  \texttt{\{yb2630, kmb2297, wk2439\}@columbia.edu}
}

\begin{document}

\maketitle

\begin{abstract}
Wish-driven programming---stating an intent and letting an agent build the artifact---fails
quietly: the program runs and still violates unwritten requirements. We delete fields
from full specifications mechanically and train a small \emph{silence detector} to predict which
requirement categories a wish leaves implicit, then target clarification questions. This
milestone reports committed artifacts only: 1{,}100 kept tasks, a five-fold out-of-fold detector,
a 45-unit dry run, and a prepared but unrun 100-task pilot. The encoder reaches macro-F1 0.777
(0.629 varying) and AUROC 0.778/0.947/0.736 for examples/constraints/error behaviour; TF-IDF
reaches 0.782/0.909/0.959/0.734, so the encoder is not the strongest detector. No
condition-level results exist yet.
\end{abstract}

\section{Approach}

\paragraph{Task and unit.}
The agent receives a one-sentence wish and a Python stub and must pass a hidden unit-test suite.
A unit is one function with a shared stub, visible tests, and hidden suite. We use HumanEval+ and
MBPP+ \citep{liu2023evalplus} plus BigCodeBench; the first two pair a docstring with one function,
and BigCodeBench is our library-call stress case.

\paragraph{Levels by field deletion.}
Each specification is parsed into intent, signature and types, examples, constraints and edge
cases, and error behaviour. L0 is full; L1 drops examples and type detail; L2 additionally drops
constraints and error behaviour; L3 keeps only intent. Deleted fields are the detector labels,
and single-field drops L2$_c$ remove one category from L0. The design uses the \emph{task
sentence} as intent, not the first sentence: HumanEval/104's first sentence only says ``Given a
list of positive integers x.'', while its task sentence states the requirement. The committed R0
artifact still applies a first-sentence rule, so HumanEval/104's L3 is that precondition; the
prepared pilot uses this known-imperfect wish, and a task-sentence extractor is future work.

\paragraph{Visible tests.}
Visible tests come from the hidden suite and are fixed across levels. The design wants them
disjoint from the examples; the artifact does not. Of 2{,}608 visible calls with
literal arguments, 1{,}311 (50.3\%) also occur in the examples, and MBPP+ contributes 1{,}131 of
1{,}131 because its examples field is the test list. On the disjoint 100-task pilot control the
overlap is 121 of 232 (52.2\%). We report the measurement instead of asserting disjointness.

\paragraph{Harness and detector.}
The agent is a ReAct loop with write, run-visible-tests, read-error and ask-stakeholder tools, a
fixed turn cap, and no repository; H4 gives sibling functions. The stakeholder sees L0 and
answers only from it \citep{mu2023clarifygpt}. Contemporary harnesses revise state
online \citep{xiao2026reharness}; we hold harness state fixed. The detector is a frozen
Qwen3-Embedding-0.6B encoder with a multi-label head, BCE and an L2 penalty, evaluated five-fold
out-of-fold with inner-fold calibration. A positive label means the category has
L0 content the wish omits. The H3 policy ranks categories by calibrated presence times utility, asks
2.47 questions per task on average, and never asks intent; answers become a Requirements block
and a defensive-coding instruction.

\paragraph{Conditions and related work.}
H0 is the plain loop; H1 writes acceptance tests first; H2 asks generic questions; H3 asks
detector-targeted questions; H4 exposes sibling functions; H5, planned, evolves the harness
offline on 200 development wishes with RRSI's edits \citep{xia2026rrsi}. B1 gives L0 as the
ceiling; B2 spends the budget on extra samples; B3 is a frontier model given only the wish;
DISCARD asks H3's questions and withholds answers. The prepared pilot runs H0--H4, B1--B3, and
DISCARD. Harness$^2$ evolves harnesses without ground-truth verifiers, reporting up to 11.5
points on WorkBuddy-Bench, 4.5 on JobBench, and 11.8 percentage points on LAB
\citep{xu2026harness2}. HarnessPro evolves harnesses label-free with LLM-generated executable
tests instead of an LLM judge, reporting +2.0 percentage points and 26.1\% lower per-task cost
than RHO in its heavy mode \citep{pan2026harnesspro}. H5 is verifier-free, so these works bound it.

\section{Experiments}

\paragraph{Data and splits.}
Of 1{,}682 source rows, 1{,}676 parse and 1{,}573 validate; 1{,}100 are kept (157 HumanEval+,
377 MBPP+, 566 BigCodeBench), with 109 exclusions. The split is 500 train, 200 evolution, 300
held-out, and a disjoint 100-task pilot (seed 20261121). Kept-set prevalence is 1.000 (intent,
signature/types), 0.973 (examples), 0.645 (constraints), and 0.134 (error behaviour); error
behaviour has 63 train and 147 kept positives. Held-out tasks are untouched.

\paragraph{Detector results.}
Table~\ref{tab:detector} gives five-fold out-of-fold detection. The
encoder reaches macro-F1 0.777 (0.629 varying) and AUROC 0.778/0.947/0.736 for examples,
constraints, and error behaviour. TF-IDF reaches 0.782/0.909/0.959/0.734, and the prior baseline 0.753/0.588. The encoder is 0.004/0.131/0.012 below TF-IDF on
macro-F1/examples/constraints and slightly above on error behaviour (0.736 vs 0.734);
calibration is comparable (Brier 0.042 vs 0.041, ECE 0.012 vs 0.013). We treat this gap as a
limitation.

\begin{table}[!ht]
  \centering
  \footnotesize
  \setlength{\tabcolsep}{4pt}
  \begin{tabular}{@{}lccccc@{}}
    \toprule
    Method & macro-F1 & macro-F1 (varying) & AUROC examples & AUROC constraints & AUROC error \\
    \midrule
    Prior & 0.753 & 0.588 & 0.358 & 0.426 & 0.409 \\
    TF-IDF & 0.782 & 0.636 & 0.909 & 0.959 & 0.734 \\
    Encoder & 0.777 & 0.629 & 0.778 & 0.947 & 0.736 \\
    \bottomrule
  \end{tabular}
  \caption{Five-fold out-of-fold detection on the 500 training tasks. ``Varying'' averages
  examples, constraints and error behaviour; intent and signature/types are positive on every
  task, so their AUROC is undefined. Source: \texttt{results/detector/metrics.json}.}
  \label{tab:detector}
\end{table}

\paragraph{Pilot status.}
The pilot is prepared as two config-stable shards and has not run. A 45-unit dry run (five tasks, all
conditions) completed on one A100 with 886{,}758 tokens and 408\,s; all units completed but
only 1 passed, so it is a plumbing check, not a condition result. No condition-level results
exist.

\section{Future Work}

We will run the prepared pilot, then scale to the 300 held-out tasks and more seeds. First,
replace the first-sentence L3 with a task-sentence extractor and make visible tests disjoint from
examples; MBPP+ duplicates its test list in both fields. We will run H5, finish DISCARD, compare
the encoder with TF-IDF, and add a source with explicit error-behaviour annotations. The
proposal-feedback changes are in place: mechanical field deletion replaced LLM rewriting, and
the harness question is the H4 experiment; the pilot will settle it.

\bibliographystyle{unsrt}
\bibliography{references}

\end{document}
